% !TEX root = ../main.tex
% Figure content only
\begin{tikzpicture}[x=1cm,y=1cm]
  \draw[->,thick] (-3.8,-1.8) -- (3.8,-1.8) node[right] {$z$};
  \draw[->,thick] (-3.8,-1.8) -- (-3.8,2.6) node[above] {能量};

  % Quantum well region
  \fill[blue!8] (-1.2,-1.8) rectangle (1.2,2.4);
  \node[blue!60!black,font=\footnotesize] at (0,2.25) {量子阱区};

  % Bands
  \draw[thick,blue] (-3.2,1.9) -- (-1.2,1.9) -- (-1.2,1.2) -- (1.2,1.2) -- (1.2,1.9) -- (3.2,1.9)
    node[right] {$E_c$};
  \draw[thick,red] (-3.2,-0.1) -- (-1.2,-0.1) -- (-1.2,-0.7) -- (1.2,-0.7) -- (1.2,-0.1) -- (3.2,-0.1)
    node[right] {$E_{\mathrm{HH}}$};
  \draw[thick,green!50!black] (-3.2,-0.45) -- (-1.2,-0.45) -- (-1.2,-1.0) -- (1.2,-1.0) -- (1.2,-0.45) -- (3.2,-0.45)
    node[right] {$E_{\mathrm{LH}}$};

  % Quantized levels
  \draw[dashed,blue] (-1.05,1.45) -- (1.05,1.45) node[right,font=\footnotesize] {$E_{c1}$};
  \draw[dashed,red] (-1.05,-0.45) -- (1.05,-0.45) node[right,font=\footnotesize] {$E_{hh1}$};
  \draw[dashed,green!50!black] (-1.05,-0.82) -- (1.05,-0.82) node[right,font=\footnotesize] {$E_{lh1}$};

  % Transition
  \draw[->,ultra thick,orange] (0,1.4) -- (0,-0.42);
  \node[right,orange!80!black,font=\footnotesize] at (0.05,0.55) {$h\nu$};

  % Split annotation
  \draw[<->,thick] (2.7,-0.82) -- (2.7,-0.45);
  \node[right,font=\footnotesize] at (2.7,-0.64) {应变分裂};

  % Fermi levels
  \draw[dash dot,thick,purple] (-3.1,1.05) -- (3.1,1.05) node[right] {$E_{Fc}$};
  \draw[dash dot,thick,brown] (-3.1,-1.18) -- (3.1,-1.18) node[right] {$E_{Fv}$};

  % Notes
  \node[draw=black,rounded corners,fill=gray!6,align=left,font=\footnotesize,text width=6.8cm] at (0,-2.95) {
    压缩应变量子阱会把重空穴带和轻空穴带拉开，
    从而改变 TE/TM 增益排序。
    教学上最关心的是导带基态与重空穴基态之间的主跃迁。
  };
\end{tikzpicture}
